\documentclass[11pt,a4paper]{article}
\usepackage[margin=1in]{geometry}
\usepackage[T1]{fontenc}
\usepackage[utf8]{inputenc}
\usepackage{lmodern}
\usepackage{amsmath,amssymb,amsthm}
\usepackage{hyperref}

\newtheorem{proposition}{Proposition}
\newtheorem{exercise}{Exercise}

\title{Beginner Guide to Informed Machine Learning\\from the Papers in This Folder}
\author{}
\date{\today}

\begin{document}
\maketitle

\section{How to Read This Guide}
This guide assumes you are a beginner. We use one principle:
\begin{quote}
\textbf{First intuition, then formula, then what to prove yourself.}
\end{quote}

\section{Big Picture in One Sentence}
Classical machine learning says: ``learn from data.''\newline
Informed machine learning says: ``learn from data \emph{and} known knowledge (rules, equations, physics, expert logic).''

\section{Part A: Taxonomy (What is Being Added to ML?)}
From the survey paper, every informed ML method can be understood by 3 questions:
\begin{enumerate}
\item Where does knowledge come from? (science, experts, world knowledge)
\item How is knowledge represented? (logic rules, equations, simulations, graphs)
\item Where is it integrated? (data, model structure, training loss, final output)
\end{enumerate}

\subsection{Simple mental model}
Think of normal ML as cooking with only ingredients (data).\newline
Informed ML adds recipe constraints (knowledge).

\section{Part B: Three Logic Papers (How to Teach a Neural Network Rules)}

\subsection{1) Logic-net (teacher-student idea)}
\textbf{Intuition:} Build a ``teacher'' that obeys rules more, then make your network imitate that teacher.

\textbf{Core form:}
\begin{equation}
\mathcal{L}(\theta)= (1-\pi)\,\mathcal{L}_{\text{label}} + \pi\,\mathcal{L}_{\text{teacher}}.
\end{equation}

So the student learns from true labels and rule-aware teacher outputs.

\subsection{2) Semantic Loss (probability of satisfying logic)}
\textbf{Intuition:} If your output should satisfy a logic statement $\alpha$, reward the model when probability mass is on worlds that satisfy $\alpha$.

\textbf{Core form:}
\begin{equation}
\mathcal{L}_{s}(\alpha,p)
= -\log \sum_{x \models \alpha}
\prod_{i: x_i=1} p_i \prod_{i: x_i=0}(1-p_i).
\end{equation}

If satisfying worlds have high probability, this loss is small.

\subsection{3) DL2 (logic language to differentiable loss)}
\textbf{Intuition:} Write constraints declaratively, automatically translate them into a smooth loss, then optimize.

\textbf{Training form:}
\begin{equation}
\min_{\theta}\;\mathcal{L}_{\text{task}}(\theta)+\lambda\,\mathcal{L}_{\text{logic}}(\theta).
\end{equation}

\textbf{Query form:}
\begin{equation}
\min_{x}\;\mathcal{L}_{\text{query-constraint}}(x,\theta)
\end{equation}
which means: find an input that satisfies desired logic.

\section{Part C: Granular Computing Papers (How to Express Confidence)}

\subsection{Why needed?}
A single number prediction can look precise but hide uncertainty.

\subsection{Core idea}
Instead of outputting one value $\hat y$, output a granule:
\begin{itemize}
\item interval: $[l,u]$
\item fuzzy set: graded confidence
\item probabilistic granule: e.g., Gaussian $\mathcal{N}(\mu,\sigma^2)$
\end{itemize}

\subsection{Takagi-Sugeno baseline}
\begin{equation}
\hat y(x)=\frac{\sum_{i=1}^{c} w_i(x)b_i}{\sum_{i=1}^{c} w_i(x)}.
\end{equation}

Granular version replaces scalar $b_i$ by interval/fuzzy/probabilistic consequents.

\subsection{Justifiable Granularity principle}
Choose a granule that balances coverage and specificity:
\begin{equation}
G^*=\arg\max_G\;V(G),\qquad V(G)=\mathrm{cov}(G)\cdot\mathrm{sp}(G).
\end{equation}

\subsection{GP bridge}
If prediction is Gaussian:
\begin{equation}
y(x)\mid\mathcal D\sim \mathcal N(\mu(x),\sigma^2(x)),
\end{equation}
then a confidence interval is
\begin{equation}
[\mu-z\sigma,\;\mu+z\sigma].
\end{equation}

\section{Part D: Unified Formula You Should Remember}
\begin{equation}
\min_{\theta}\;\mathcal L_{\text{task}} + \lambda_1\mathcal L_{\text{knowledge}} + \lambda_2\mathcal L_{\text{uncertainty}}.
\end{equation}

This single equation unifies most papers in your folder.

\section{Research Propositions for You to Prove}
Below are beginner-friendly propositions/exercises. Try proving them step by step.

\begin{proposition}[Monotonicity of semantic loss contribution]
Fix a constraint $\alpha$ and probabilities of all worlds except one satisfying world $x^*$. If the model increases probability mass assigned to satisfying worlds (without violating normalization), then semantic loss does not increase.
\end{proposition}
\textit{Hint:} semantic loss is $-\log(\text{satisfying mass})$.

\begin{exercise}[Proof task]
Formalize the proposition above by defining $S=\sum_{x\models\alpha}P_p(x)$ and show that if $S'\ge S$, then $-\log S'\le -\log S$.
\end{exercise}

\begin{proposition}[Teacher-student interpolation limits]
For loss
\(\mathcal L=(1-\pi)\mathcal L_{\text{label}}+\pi\mathcal L_{\text{teacher}}\),
when $\pi=0$ it reduces to pure supervised learning; when $\pi=1$ it reduces to pure teacher imitation.
\end{proposition}

\begin{exercise}[Proof task]
Write the two substitutions explicitly and interpret what signal disappears in each case.
\end{exercise}

\begin{proposition}[Coverage-specificity tradeoff]
Suppose specificity is strictly decreasing with interval width $w$ and coverage is nondecreasing with $w$. Then maximizing
\(V(w)=\mathrm{cov}(w)\mathrm{sp}(w)\)
produces an interior optimum under mild smoothness/concavity assumptions.
\end{proposition}

\begin{exercise}[Proof task]
Assume differentiability and solve $\frac{d}{dw}\log V(w)=0$, i.e.
\[
\frac{\mathrm{cov}'(w)}{\mathrm{cov}(w)} = -\frac{\mathrm{sp}'(w)}{\mathrm{sp}(w)}.
\]
Interpret this as balancing marginal gain and marginal loss.
\end{exercise}

\begin{proposition}[Interval output contains point output]
If each rule consequent interval collapses to a point, i.e. $[b_i,b_i]$, then interval aggregation returns
\([\hat y,\hat y]\)
where $\hat y$ is the classical TS output.
\end{proposition}

\begin{exercise}[Proof task]
Substitute $\underline b_i=\overline b_i=b_i$ into interval aggregation and simplify both bounds.
\end{exercise}

\begin{proposition}[Calibration motivation]
For Gaussian predictive distribution, fixed confidence level $1-\alpha$ gives interval width proportional to $\sigma(x)$.
Hence higher uncertainty points naturally get wider intervals.
\end{proposition}

\begin{exercise}[Proof task]
From $[\mu-z_{\alpha/2}\sigma,\mu+z_{\alpha/2}\sigma]$, compute width and show linear dependence on $\sigma$.
\end{exercise}

\section{Suggested Study Path (7 Days)}
\begin{enumerate}
\item Day 1--2: Read taxonomy paper; map each other paper to source/representation/integration.
\item Day 3: Re-derive semantic loss on a tiny 2-variable logic constraint.
\item Day 4: Re-derive teacher-student objective and test extreme $\pi$ values.
\item Day 5: Implement one toy DL2-style constraint loss.
\item Day 6: Build interval output from a simple TS model.
\item Day 7: Prove two propositions above and write a short report.
\end{enumerate}

\section{Primary References}
\begin{itemize}
\item L. von Rueden et al., IEEE TKDE, 2021/2023 version.
\item Z. Hu et al., \emph{Harnessing Deep Neural Networks with Logic Rules}.
\item J. Xu et al., \emph{A Semantic Loss Function for Deep Learning with Symbolic Knowledge}.
\item M. Fischer et al., \emph{DL2: Training and Querying Neural Networks with Logic}.
\item Y. Cui et al., \emph{From Fuzzy Rule-Based Models to Granular Models}, IEEE TFS 2025.
\item W. Pedrycz, \emph{Granular Computing for Machine Learning}, IEEE TCYB 2025.
\end{itemize}

\end{document}
